\chapter{Robustness \& Diagnostics}
\label{chap:robustness}

The results chapter ranked the channels by how clearly and durably they respond to an ECB easing surprise. This chapter asks whether that ranking survives when the specification, the shock, and the sample are changed. Here robustness has the same meaning it has in the research question. It is durability of the sign and of the accumulated direction, not identical coefficients in every cut of the data. A response fails the test when either one flips. The point of the exercise is to see whether the comparative ordering in \Cref{tab:persistence-synthesis} is a feature of the data or an artifact of one specification.

Four responses cluster at the stable end of the checks below. Mortgage spreads, NFC spreads, house-purchase lending growth, and the real DAX hold their direction across the widest set of checks. These are the intermediation and housing-credit variables that carry the mechanism, plus the equity price that sits alongside it. The compensation-linked proxies stay informative but move more with the specification.

\section{Alternative Lag Structures}

The pipeline does not run an exhaustive lag grid, and this is a design choice rather than an omission. The horizons are the irregular set \(\{0,1,3,6,12,24\}\), so the natural threat is not a mis-specified lag order but the overlap between neighbouring projections. That overlap is handled by the HAC and bootstrap machinery set out in \Cref{chap:methodology}, together with the rolling and recursive windows re-estimated below.

Two layers carry that stress. The seven event and exclusion samples defined in \Cref{sec:subsample-analysis} re-estimate every canonical response at every horizon, and rolling and recursive windows re-estimate four of them. The rolling windows are 84 months long and move one month at a time. The recursive estimates start from the first 60 months and extend in six-month steps. Both run for house-purchase lending growth, pure-new loans, the wage tracker, and the DAX at the six- and twelve-month horizons. Across the sample cuts the same four responses stay put. House-purchase lending growth, mortgage spreads, NFC spreads, and the DAX keep their modal sign in at least 86 percent of the sample-by-horizon cells and their accumulated direction in every sample. NFC credit and the wage tracker do not, and their signs move once the event sample or the exclusion window shifts. The rolling windows are the harsher test. House-purchase lending growth keeps its positive sign in about three quarters of the 84-month windows at six months and two thirds at twelve, and in every recursive sample at six months. The wage tracker is negative in almost every six-month window. The DAX sign at six months is close to a coin toss across windows, while at twelve months it is negative in two thirds of them, which is why the equity result is read as a medium- and long-horizon result.

The channels that answer the research question most clearly are also the ones whose direction is least sensitive to how the dynamics are specified.

\section{Alternative Shock Measures}

The baseline shock is the easing-signed target surprise. The timing, forward-guidance, QE, and weighted-composite factors are re-estimated as alternative shocks to see how far the ordering depends on which piece of the surprise is used. The weighted composite sits closest to baseline. The DAX stays negative, house-purchase lending growth stays positive at h6 and h12, and the mortgage spread accumulates positive. The NFC spread is the one place the weighted composite departs from baseline, with an accumulated path that stays slightly negative even though its marginal response turns positive at the longer horizons. The forward-guidance, timing, and QE components change several signs, most visibly for wage pressure and for some housing-credit specifications. Component choice is material, which is the expected reading once the surprise is decomposed rather than treated as one object \parencite{JarocinskiKaradi2020}. These alternative factors serve as diagnostics on the shock and leave the baseline in place. The baseline and alternative-shock estimates share the normalized IRF output reported in \Cref{tab:normalized-irf-output}.

A separate check compares the surprise-based local projection against a plain OLS regression on the raw policy-rate change. The raw-rate regression is estimated for the ten canonical responses, which gives 60 comparable response-horizon cells. Read in the same direction, with a rate cut counted as easing, the two designs agree in 36 of those cells and disagree in 24. Full agreement is not expected. The surprise design uses easing-signed announcement news, while the raw policy-rate regression uses the realised rate move, so the two objects should not track each other coefficient by coefficient. Every cell is labelled as a qualitative direction check, and the file treats it as a sign comparison with no causal reading. The file compares raw signs, so its match column records the opposite of the direction-adjusted reading used here. The lending spreads and the wage tracker agree at every horizon, and the DAX through twelve months. House-purchase lending growth agrees at six and twelve months, while employment expectations and the German wage bill disagree at nearly every horizon.

\section{Sub-Sample Analysis}
\label{sec:subsample-analysis}

The sample is cut seven ways: the full sample, clean events only and contaminated events only, which rebuild the monthly shock from the two event classes of the screen, and four historical exclusions that drop the COVID era from January 2020 to December 2021, the crisis windows from September 2008 to June 2009, July 2011 to December 2012, and March 2020 to December 2021, the five months with the largest absolute surprises, and the QE launch from January to June 2015. The cumulative results are stable where it counts. House-purchase lending growth accumulates positive at h6, h12, and h24 in all seven samples. Mortgage and NFC spreads accumulate positive at the same three horizons in all seven. The DAX accumulates negative at h12 and h24 in all seven.

The fragile responses behave differently under the same cuts. The wage tracker changes accumulated sign between samples. It is positive in the clean-event sample and negative in the full and contaminated samples, so its long-run direction is not a fixed feature of the data. NFC credit is inconsistent across the seven cuts and never settles on one direction. These are the same two responses that moved under the lag and window checks.

One distinction matters for how the house-purchase result is read. The durable positive claim is about the accumulated path, not the marginal coefficient at every horizon. The marginal response for house-purchase lending growth is positive through h12 and turns negative at h24 in six of the seven filters, with only the crisis-window exclusion keeping it positive. The accumulated path still stays positive because the earlier horizons carry more weight. This is the same object seen in the baseline, where the h24 point coefficient is negative while the running sum remains positive. The directional matrix records both the marginal and the cumulative signs, which is why the two are kept separate.

\thesisfigure
  {enhancement/robustness_consistency_matrix.png}
  {Robustness consistency matrix for cumulative response directions across event filters and historical exclusions.}
  {fig:robustness-consistency-matrix}

\Cref{fig:robustness-consistency-matrix} shows why the claim is stated in accumulated terms. Housing-credit growth and the lending spreads keep their cumulative direction across the main filters at medium and long horizons. The wage tracker and employment expectations move with sample composition.

\section{Identification Fragility}

Not every announcement is a clean monetary signal, and the event screen quantifies how much of the sample is exposed. Of 244 events, 86 look like possible pure monetary shocks (35.2 percent), 81 show a mixed sign conflict between the rate factor and the yield response (33.2 percent), 68 look like possible information shocks (27.9 percent), and 9 are ambiguous zero responses (3.7 percent). Roughly a third of the events read as clean surprises and the rest carry some contamination or conflict. The share shifts by regime. The financial-crisis window has the highest share of possible information shocks, close to half of its events. These counts describe the full 244-event archive. The estimation samples use the narrower monthly identification window from \Cref{chap:methodology}, so the clean-event count there is a subset of this archive total.

\begin{table}[htbp]
  \centering
  \caption{Information-effect classification of ECB announcement events}
  \label{tab:information-effect-screen}
  \footnotesize
  \setstretch{1.0}
  \setlength{\tabcolsep}{6pt}
  \begin{tabular}{@{}l S[table-format=3.0] S[table-format=3.1]@{}}
    \toprule
    Event classification & \multicolumn{1}{c}{Events} & \multicolumn{1}{c}{Share (\%)} \\
    \midrule
    possible pure monetary shock & 86 & 35.2 \\
    mixed sign factor yield conflict & 81 & 33.2 \\
    possible information shock & 68 & 27.9 \\
    ambiguous zero response & 9 & 3.7 \\
    \midrule
    Total & 244 & 100.0 \\
    \bottomrule
  \end{tabular}
  \vspace{0.3em}
  \begin{minipage}{\textwidth}\footnotesize\setstretch{1.0}
  Notes: Events are classified by comparing the sign of the rate factor with the
  yield and equity responses in the announcement window. Counts cover the full
  event archive.
  Source: own calculations.
  \end{minipage}
\end{table}


The screen sorts the responses the same way the other checks do. House-purchase lending growth and the lending spreads keep their qualitative direction across clean and contaminated events. The DAX is more negative in the contaminated events, which is the pattern expected if easing announcements also carry adverse macroeconomic news. Because the screen itself uses the sign of the equity move, part of this difference is built into the classification, so the information-effect reading remains an interpretation rather than a tested result. The wage tracker changes accumulated sign between the clean and contaminated samples. So the housing-credit and lending-margin ordering survives the contamination screen, while the equity and wage readings are the ones that need the information-effect caveat attached.

\section{Relevance Sensitivity}

The strongest first-stage bridge is the end-of-month DFR change on the easing-signed target surprise, with \(F=17.07\) and partial \(R^2=0.066\). The full-sample first stage clears the weak-instrument bar on the classical statistic. The problem is where the relevance comes from. Drop the five largest absolute surprises, which make up 28.8 percent of the total absolute shock, and the first-stage \(F\) falls to 5.34, about 0.31 of its full-sample value. The rolling median \(F\) is 8.99, below the Stock-Yogo-style thresholds, even though the sign of the first-stage coefficient never flips across the rolling windows. Relevance is real but concentrated in a handful of large events.

\begin{table}[htbp]
  \centering
  \caption{Monthly first-stage relevance of the strongest candidate policy bridges}
  \label{tab:instrument-strength-matrix}
  \footnotesize
  \setstretch{1.0}
  \setlength{\tabcolsep}{5pt}
  \begin{tabular}{@{}ll S[table-format=2.2] S[table-format=1.3] S[table-format=3.0] l@{}}
    \toprule
    Surprise factor & Policy target & \multicolumn{1}{c}{$F$} & \multicolumn{1}{c}{Partial $R^2$} & \multicolumn{1}{c}{$N$} & Screen \\
    \midrule
    target factor & deposit facility rate, end of month & 17.07 & 0.066 & 246 & fragile candidate \\
    timing factor & ECB assets, log monthly average & 10.99 & 0.044 & 244 & fragile candidate \\
    fg factor & deposit facility rate, monthly average & 9.43 & 0.038 & 244 & fragile candidate \\
    timing factor & ECB assets, log end-of-month stock & 7.05 & 0.029 & 244 & fragile candidate \\
    qe factor & deposit facility rate, end of month & 5.73 & 0.040 & 142 & fragile candidate \\
    target factor & ECB assets, log monthly average & 4.42 & 0.018 & 246 & reject weak relevance \\
    target factor & household credit & 3.02 & 0.012 & 246 & reject weak relevance \\
    qe factor & NFC credit & 2.50 & 0.018 & 142 & reject weak relevance \\
    \bottomrule
  \end{tabular}
  \vspace{0.3em}
  \begin{minipage}{\textwidth}\footnotesize\setstretch{1.0}
  Notes: Eight strongest bridges from the monthly first-stage tournament, ranked by
  $F$-statistic. No candidate clears the strength and stability screen jointly, which
  is why the thesis reports reduced-form responses to policy news.
  Source: own calculations.
  \end{minipage}
\end{table}


The tournament makes the same point across the full field. None of the 45 candidate bridges passes cleanly. The five strongest are flagged as fragile candidates and the remaining forty are rejected for weak relevance. A full-sample first stage is informative, but it is not steady enough to support a strong claim about one treatment variable when the relevance is event-concentrated and regime-sensitive \parencite{StockYogo2005}. The thesis takes the conservative route. The surprise is not upgraded into an instrument for the DFR, the shadow rate, ECB assets, or the credit aggregates, and the estimand stays the reduced-form policy-news response. This narrows the treatment interpretation. It does not weaken the persistence comparison, which never depended on a structural first stage.

\section{Sign Consistency}

The sign-consistency ranking lines up with everything above and with the stability scores already tabulated in \Cref{tab:persistence-synthesis}. The lending spreads and house-purchase lending growth sit at the top, the DAX matches them in strength but in the opposite direction, and NFC credit is the weakest, close to a coin toss and too imprecise to carry a persistence claim.

\thesisfigure
  {uncertainty/significance_heatmap.png}
  {Horizon-by-response significance heatmap for thesis-facing response variables.}
  {fig:significance-heatmap}

Individual coefficient significance is thin, and this is why the argument is built on sign patterns rather than single coefficients. Most horizon-by-response cells in the heatmap are not significant on their own. The clearest cases include house-purchase lending growth at the short horizons, the DAX at h12 and h24, the German wage bill at h0 and h6, employment expectations at h24, the mortgage spread at h1, and the NFC spread at h12. The persistence claim rests on the accumulated bands, where the cumulative confidence pattern is stronger. Under the HAC intervals, cumulative significance holds across the whole horizon path only for house-purchase lending growth and the German wage bill. It holds from h1 for the mortgage spread and from h6 for the NFC spread. For the DAX it holds only at h12 and h24. NFC credit, ECB assets, pure-new loans, household credit, and retail volume never reach it.

The long-horizon inference is joint rather than coefficient-by-coefficient because the projections overlap. Adjacent-horizon residual correlations run from 0.41 to 0.82 and reach 0.87 for NFC credit and 0.86 for house-purchase growth. That dependence is the reason the uncertainty layer uses moving-block and wild bootstraps and reads the cumulative bands as a set. The moving-block bootstrap uses 399 replications with a block length of 12, and the wild bootstrap gives a second view on the same responses.

\Cref{tab:residual-diagnostics} reports four tests on the residuals of every baseline regression, so that the choice between these interval methods rests on evidence rather than on convention. The first asks whether serial correlation remains beyond the HAC bandwidth, using the twelve autocorrelations at lags \(h+2\) to \(h+13\). It rejects at the five percent level in 45 of the 60 cells, and at every horizon for ECB assets, NFC credit, house-purchase lending growth, and pure new loans. The bandwidth \(h+1\) therefore covers the overlap that the design builds in, but not all of the dependence that remains, which is why the HAC intervals are reported alongside, and never instead of, the joint bootstraps, which resample all horizons together. The second and third tests concern the variance. Koenker's form of the Breusch--Pagan test rejects constant variance in 34 cells, and Engle's ARCH test with twelve lags rejects in 51, so the residual variance moves over time for almost every response. That is the condition the wild bootstrap is built for, and it is the reason the wild interval is reported next to the moving-block interval rather than treated as a secondary check. The fourth test, Jarque--Bera, rejects normality in 48 cells, with excess kurtosis the cause in 47 of them, which is why every bootstrap interval is a percentile interval and none relies on a normal approximation.

\begin{table}[htbp]
  \centering
  \caption{Residual diagnostics for the baseline local projections}
  \label{tab:residual-diagnostics}
  \footnotesize
  \setstretch{1.0}
  \setlength{\tabcolsep}{4pt}
  \begin{tabular}{@{}l r l l l l@{}}
    \toprule
    & & \multicolumn{4}{c}{Horizons (months) at which the test rejects at 5 percent} \\
    \cmidrule(lr){3-6}
    Response & \multicolumn{1}{c}{$N$} & Serial corr. & Heterosk. & ARCH & Non-normal \\
    \midrule
    house-purchase lending growth & 246 & all & 1, 3, 6, 12, 24 & all & 0, 1, 3 \\
    pure new house-purchase loans & 97 & all & 3, 6 & 12 & none \\
    mortgage lending spread & 246 & 6, 12, 24 & 1, 3, 6, 12, 24 & all & all \\
    NFC lending spread & 246 & 0, 12, 24 & 1 & all & all \\
    NFC credit & 246 & all & 6, 12, 24 & 1, 3, 6, 12, 24 & 0, 1, 3, 6, 12 \\
    real DAX & 246 & 6, 12, 24 & all & 1, 3, 6, 12, 24 & all \\
    ECB assets & 246 & all & 0, 1, 3 & all & 0, 1, 3, 24 \\
    real wage tracker excl. one-offs & 152 & 0, 1, 3, 6, 12 & 0, 1, 6, 12, 24 & 1, 3, 6, 12, 24 & all \\
    German industry real wage bill growth & 235 & 0, 1, 3, 24 & 1, 3 & all & all \\
    employment expectations & 246 & 3, 6, 12 & 0, 1 & 1, 3, 6, 12, 24 & all \\
    \bottomrule
  \end{tabular}
  \vspace{0.3em}
  \begin{minipage}{\textwidth}\footnotesize\setstretch{1.0}
  Notes: Baseline target-shock regressions, full sample. $N$ is the estimation sample at $h=0$.
  Serial corr.: Ljung--Box test on the twelve residual autocorrelations at lags $h+2$ to $h+13$,
  just beyond the Newey--West bandwidth $h+1$. Heterosk.: Koenker's studentized Breusch--Pagan
  test of the squared residuals on the regressors. ARCH: Engle's LM test with twelve lags.
  Non-normal: Jarque--Bera. ``all'' means every one of the six horizons. Source: own calculations.
  \end{minipage}
\end{table}


Read together, the diagnostics overturn no result, but they fix the order in which the intervals should be trusted at the long horizons. Where the HAC interval and the joint bootstraps agree, as they do for housing credit through twelve months, the conclusion is firm. Where they disagree, as they do for housing credit at twenty-four months, the bootstrap reading takes precedence, and \Cref{chap:results} is written that way.

\begin{table}[htbp]
  \centering
  \caption{Accumulated responses: 90 percent intervals from the dependence-aware bootstraps}
  \label{tab:bootstrap-cumulative}
  \footnotesize
  \setstretch{1.0}
  \setlength{\tabcolsep}{4pt}
  \begin{tabular}{@{}llccc@{}}
    \toprule
    & & \multicolumn{3}{c}{Horizon in months} \\
    \cmidrule(lr){3-5}
    Response & Bootstrap & 6 & 12 & 24 \\
    \midrule
    house-purchase lending growth & moving block & [0.167, 0.545]$^{\dagger}$ & [0.109, 0.731]$^{\dagger}$ & [-0.078, 0.734] \\
     & wild & [0.143, 0.644]$^{\dagger}$ & [0.151, 0.819]$^{\dagger}$ & [-0.063, 0.809] \\[0.3em]
    pure new house-purchase loans & moving block & [-0.042, 0.082] & [-0.065, 0.091] & [-0.088, 0.073] \\
     & wild & [-0.059, 0.101] & [-0.068, 0.107] & [-0.085, 0.063] \\[0.3em]
    mortgage lending spread & moving block & [0.016, 0.134]$^{\dagger}$ & [0.039, 0.218]$^{\dagger}$ & [0.034, 0.267]$^{\dagger}$ \\
     & wild & [-0.009, 0.148] & [-0.001, 0.234] & [-0.016, 0.293] \\[0.3em]
    NFC lending spread & moving block & [-0.004, 0.081] & [0.008, 0.125]$^{\dagger}$ & [-0.011, 0.155] \\
     & wild & [-0.031, 0.106] & [-0.024, 0.149] & [-0.040, 0.176] \\[0.3em]
    NFC credit & moving block & [-0.005, 0.003] & [-0.008, 0.006] & [-0.013, 0.009] \\
     & wild & [-0.008, 0.006] & [-0.012, 0.010] & [-0.018, 0.013] \\[0.3em]
    real DAX & moving block & [-0.044, 0.011] & [-0.075, 0.000] & [-0.111, -0.022]$^{\dagger}$ \\
     & wild & [-0.044, 0.010] & [-0.074, -0.006]$^{\dagger}$ & [-0.123, -0.011]$^{\dagger}$ \\[0.3em]
    ECB assets & moving block & [-0.035, 0.015] & [-0.058, 0.021] & [-0.084, 0.026] \\
     & wild & [-0.049, 0.028] & [-0.078, 0.040] & [-0.114, 0.058] \\[0.3em]
    real wage tracker excl. one-offs & moving block & [-0.530, 0.061] & [-0.726, 0.175] & [-0.578, 0.533] \\
     & wild & [-0.475, -0.014]$^{\dagger}$ & [-0.572, -0.001]$^{\dagger}$ & [-0.362, 0.332] \\[0.3em]
    German industry real wage bill growth & moving block & [0.626, 2.136]$^{\dagger}$ & [0.721, 2.489]$^{\dagger}$ & [0.536, 2.458]$^{\dagger}$ \\
     & wild & [0.403, 2.356]$^{\dagger}$ & [0.404, 2.730]$^{\dagger}$ & [0.246, 2.691]$^{\dagger}$ \\[0.3em]
    employment expectations & moving block & [1.018, 4.632]$^{\dagger}$ & [0.579, 5.403]$^{\dagger}$ & [-0.462, 4.863] \\
     & wild & [-0.365, 6.132] & [-0.810, 7.365] & [-1.582, 6.679] \\[0.3em]
    \bottomrule
  \end{tabular}
  \vspace{0.3em}
  \begin{minipage}{\textwidth}\footnotesize\setstretch{1.0}
  Notes: Percentile intervals for the running sum of the normalized responses to the
  target easing surprise from 399 joint replications that resample the residuals of all
  six horizons together, in moving blocks of twelve months or with wild sign flips.
  $^{\dagger}$ marks intervals that exclude zero. Compare with the HAC intervals under
  the independence approximation in \Cref{tab:cumulative-persistence}.
  Source: own calculations.
  \end{minipage}
\end{table}


Read together, these checks say the persistence claim rests on sign-consistent accumulated paths under HAC and bootstrap uncertainty, not on any one coefficient reaching significance in isolation \parencite{Jorda2005,NeweyWest1987}.
